\section*{LỜI NÓI ĐẦU}
\thispagestyle{empty}

Trong những năm gần đây, công nghệ thiết bị bay không người lái (UAV - \mbox{Unmanned} Aerial Vehicle) đã phát triển mạnh mẽ và được ứng dụng rộng rãi trong nhiều lĩnh vực thực tiễn như tuần tra an ninh, giám sát hiện trường và trinh sát cứu nạn. Tuy nhiên, phần lớn các giải pháp hiện nay thường phụ thuộc vào các module phần cứng điều khiển thương mại nhập khẩu và phần mềm đóng gói sẵn. Việc nghiên cứu thiết kế bo mạch điều khiển từ gốc rễ và tích hợp các thuật toán xử lý trí tuệ nhân tạo (AI) thời gian thực tại trạm mặt đất vẫn là một bài toán công nghệ thách thức và mang tính thực tiễn cao đối với kỹ sư Điện tử - Viễn thông.

Đồ án \textbf{``Xây dựng và chế tạo Drone ứng dụng cho giám sát an ninh và cứu hộ''} được thực hiện nhằm xây dựng một giải pháp cứu nạn đồng bộ và tự chủ. Đồ án bao gồm hai phân hệ chính: phân hệ thiết bị bay (Drone) tập trung vào thiết kế bo mạch Flight Controller PCB (xoay quanh vi điều khiển STM32F103), tính toán phân tầng nguồn điện lực, lập trình firmware điều khiển cân bằng PID rời rạc nhúng kết hợp bộ lọc số và giải mã giao thức truyền thông vô tuyến ExpressLRS (CRSF); phân hệ trạm mặt đất (GCS) tích hợp công nghệ AI chạy song song các mô hình học sâu YOLOv8 (phát hiện đối tượng, phân loại tư thế, ước lượng 17 điểm khớp xương) kết hợp thuật toán dung hợp quyết định để đánh giá mức độ khẩn cấp (Danger Score) và tự động đưa ra các cảnh báo cứu hộ tức thời.

Dù đã có nhiều nỗ lực nghiên cứu và thực nghiệm nghiêm túc, đồ án khó tránh khỏi những thiếu sót do giới hạn về mặt thời gian và kiến thức chuyên môn. Em rất mong nhận được những ý kiến đóng góp quý báu từ các thầy, cô giáo trong Hội đồng để đồ án được hoàn thiện hơn.

Em xin bày tỏ lòng biết ơn chân thành sâu sắc tới thầy \textbf{ThS. Vũ Song Tùng}, người đã trực tiếp hướng dẫn, định hướng kỹ thuật và tận tình giúp đỡ em trong suốt quá trình hoàn thành đồ án tốt nghiệp này. Xin chân thành cảm ơn các thầy cô giáo trong Khoa, gia đình và bạn bè đã luôn động viên, hỗ trợ em trong suốt thời gian học tập và nghiên cứu.

\vspace{0.5cm}
\hfill
\begin{minipage}{0.5\textwidth}
    \centering
    \textit{Hà Nội, ngày 22 tháng 06 năm 2026} \\
    \vspace{0.2cm}
    \textbf{Người thực hiện} \\
    \vspace{2cm}
    \textbf{HOÀNG XUÂN TOÀN}
\end{minipage}
\cleardoublepage
